\section{Zielsetzung der Arbeit}
	Um die dargelegte Problemstellung zu adressieren, liegt der Fokus dieser Arbeit auf der Konzeption und Implementierung einer robusten Infrastruktur für die Anomalieerkennung innerhalb der Ladeinfrastruktur. Übergeordnetes Ziel ist die Schaffung einer stabilen Betriebsumgebung, die einen zuverlässigen, kontextbasierten Datenaustausch zwischen den einzelnen Knoten gewährleistet. 

	Ein wesentlicher Aspekt hierbei ist die Realisierung einer konsistenten und effizienten Aggregation dezentraler Informationen (siehe Kapitel~\ref{sec:generierung_der_threat_reports} (Generierung der Threat Reports)), um die Detektionsgüte der beteiligten Entitäten signifikant zu steigern. Nur durch eine performante Umgebung, welche den reibungslosen Austausch von Kontextinformationen sicherstellt, kann die in Kapitel \ref{chapter:systemarchitektur} (Systemarchitektur und Detektion Framework) detaillierte Anomalieerkennung erfolgreich implementiert und langfristig wartbar betrieben werden. Die Entwicklung dieser stabilen Kommunikations- und Verarbeitungsarchitektur bildet somit den zentralen Bestandteil der vorliegenden Arbeit.